
\begin{frame}{Appendix --- The Transfer Channel Contracts the Margin}
\hypertarget{app:transferfoil}{}
\footnotesize
In the transfer class of campaign-attack models, an attack moves a fraction
$\phi$ of the target's support away; a fraction $\alpha$ of it leaves for the
outside option instead of reaching the attacker, and the attacker loses a
fraction $\gamma$ of its own support to backlash. \ns{\citep{nakaguma2025}}
\smallskip

That model's equilibrium has every candidate target the highest-ranked
opponent, so the top two attack each other. Its evidence, measured on
electoral lawsuits, agrees: in three-candidate races $2\to1$ runs at 3.0pp,
$1\to2$ at 1.9pp and $3\to1$ at 0.8pp.
\smallskip
\begin{center}
\begin{minipage}{0.86\textwidth}
\begin{block}{Proposition 1}
\footnotesize
Under mutual $1 \leftrightarrow 2$ attacks,
\[
  s_1 - s_2 \;=\; (s^\circ_1 - s^\circ_2)\,\kappa,
  \qquad \kappa \;=\; 1 - \gamma - \phi(2-\alpha) \;<\; 1
\]
for every $\phi \in (0,1)$, $\alpha \in [0,1]$, $\gamma \in [0,1)$
satisfying the feasibility constraint $\phi + \gamma \le 1$. Adding
$3 \to 1$ subtracts a further $\phi\,s^\circ_1$. The transfer map cannot widen
the top-two margin while preserving the identity of the leader.
\end{block}
\end{minipage}
\end{center}
\smallskip

Widening requires a one-sided downward attack, which the equilibrium excludes
and the data does not show. The per-pair frequencies are also 1--3pp, too
small for the shifts the ballot decomposition reports. Something operating at
the level of the race, not between two candidates, is doing the work.
\smallskip

{\centering\scriptsize\ns{The margin is an order statistic, so a dispersion
shock widens it mechanically. The ballot and the seat split are what rule that
out.}\par}
\vfill
\centerline{\hyperlink{app:ballotdecomp}{\beamergotobutton{The ballot decomposition}}\quad\hyperlink{app:theory}{\beamerreturnbutton{Back}}}
\end{frame}
